\documentclass[11pt]{article}
\usepackage[margin=1in]{geometry}
\usepackage{amsmath,amssymb,amsthm,mathtools,hyperref}
\newtheorem{theorem}{Theorem}
\newtheorem{lemma}[theorem]{Lemma}
\newtheorem{corollary}[theorem]{Corollary}
\newcommand{\Ric}{\operatorname{Ric}}
\newcommand{\Vol}{\operatorname{Vol}}
\newcommand{\Met}{\operatorname{Met}}
\newcommand{\R}{\mathbb R}
\newcommand{\W}{\mathcal W}
\title{Euclidean-end bubbles and unrestricted curvature minimization}
\author{Research continuation}
\date{15 October 2026}
\begin{document}
\maketitle
\begin{abstract}
We give a direct application of the Chang--Gursky--Yang sphere theorem to
asymptotically Euclidean four-manifolds with Euler characteristic one.
Below a specified curvature threshold their cores are standard smooth balls.
A replacement argument then excludes positive-energy bubbles of this type
from genuine minimizing sequences for the squared traceless Ricci curvature.
This fills the trivial-end branch left open by the preceding aspherical-core
argument. It does not prove smooth attainment of the global infimum.
The sphere theorem and the analytic facts about AE ends are classical inputs;
no priority claim is made for the consequences recorded here.
\end{abstract}

\section{Conventions and the target}
For a smooth four-dimensional metric put
\[
 P(g)=\int|\Ric_g^\circ|^2\,dV_g,\qquad
 S(g)=\int s_g^2\,dV_g,\qquad
 \W(g)=\int|W_g|_{\mathrm{op}}^2\,dV_g.
\]
The tensor norm satisfies $|W|_{(0,4)}^2=4|W|_{\mathrm{op}}^2$.
On a closed manifold,
\[
 \int|\operatorname{Rm}|_{\mathrm{op}}^2=8\pi^2\chi+P.
\]
The problem under investigation asks for a closed smooth four-manifold on
which $P$ attains its unrestricted global minimum at a non-Bach-flat metric;
see LeBrun \cite[Section 11]{LeBrun}. All competitors below are smooth
metrics on the same manifold. No restriction to a conformal class is imposed.

An AE end of order $\tau>0$ is identified with $\R^4\setminus B$ and satisfies
\[
 |\partial^a(g_{ij}-\delta_{ij})|\le C_a r^{-\tau-|a|},
 \qquad |a|\le2.
\]
We assume a smooth complete connected oriented manifold with exactly one
such end and no boundary. Its compact core $N_R$ has boundary $S^3$.
The differentiated decay makes the curvature energies finite and gives
the AE Gauss--Bonnet formula
\begin{equation}\label{gb}
 8\pi^2(\chi(N)-1)=\W(g)-\tfrac12P(g)+\tfrac1{24}S(g).
\end{equation}

\section{A standard smooth core from an integral bound}
\begin{theorem}\label{standard}
Let $(N,g)$ satisfy the preceding AE assumptions, $s_g\ge0$, and
$\chi(N)=1$. If $P(g)<16\pi^2$, then the one-point compactification
defined by the AE coordinates is diffeomorphic to $S^4$.
Every sufficiently large compact core is diffeomorphic to the standard
$B^4$, with its end parametrization extendible across the core.
No Bach-flatness assumption is required.
\end{theorem}

\begin{proof}
Equation \eqref{gb} gives
\begin{equation}\label{weyl}
 \W(g)=\tfrac12P(g)-\tfrac1{24}S(g)<8\pi^2.
\end{equation}
Choose a smooth cutoff $\zeta$ equal to one for $r\le R$ and zero for
$r\ge2R$, with the usual $R^{-a}$ derivative bounds, and set on the end
\[
 g_R=\delta+\zeta(g-\delta).
\]
Extend $g_R=g$ over the core. For large $R$ these are positive definite,
uniformly equivalent to $g$, and exactly Euclidean outside $2R$.
On the transition annulus,
\[
 |\operatorname{Rm}(g_R)|\le C R^{-2-\tau},\qquad
 \Vol_{g_R}(\{R<r<2R\})\le C R^4.
\]
Consequently
\begin{equation}\label{cut}
 \|(s_{g_R})_-\|_{L^2}=O(R^{-\tau}),\qquad
 \W(g_R)\longrightarrow\W(g).
\end{equation}
The first estimate uses $s_g\ge0$ on the unchanged region.

A connected AE manifold of dimension four satisfies a homogeneous Sobolev
inequality $\|f\|_4^2\le C\int|df|^2$ for compactly supported smooth $f$.
The constant is for this fixed manifold, not for an arbitrary family.
Uniform metric equivalence gives one constant $C_0$ for all large $R$.
Thus
\[
 \int\bigl(6|df|_{g_R}^2+s_{g_R}f^2\bigr)dV_{g_R}
 \ge\bigl(6/C_0-\|(s_{g_R})_-\|_2\bigr)\|f\|_{4,g_R}^2.
\]
The coefficient is positive for large $R$. This proves positive Yamabe
constant for $g_R$. The AE Yamabe framework, including its relation to
compactification, is discussed in \cite{DM}.

Take a positive smooth function $u$ equal to $r^{-2}$ outside a still larger
compact set. The metric $\widehat g_R=u^2g_R$ extends smoothly across the
added point: in inverted coordinates it is exactly Euclidean there.
Call the resulting closed oriented manifold $\widehat N$. It has
$\chi(\widehat N)=2$. Conformal invariance of the Yamabe quotient and
the zero $W^{1,2}$ capacity of a point imply
\[
 Y(\widehat N,[\widehat g_R])=Y(N,[g_R])>0.
\]
For completeness, the equality of infima follows by cutting a smooth
test function off in a ball of radius $\varepsilon$ around the added point:
the gradient error is $O(\varepsilon^2)$ and its $L^4$ error tends to zero.
Conformal invariance of Weyl energy gives
\[
 \int_{\widehat N}|W_{\widehat g_R}|_{(0,4)}^2
 =4\W(g_R)<32\pi^2=16\pi^2\chi(\widehat N)
\]
for sufficiently large $R$, by \eqref{weyl}--\eqref{cut}.
Theorem A of Chang--Gursky--Yang \cite{CGY} now yields
$\widehat N\cong S^4$ (the $\mathbb{RP}^4$ alternative is excluded).
Their theorem gives a diffeomorphism, not conformal roundness; the former
is exactly what is needed here.

The complement of a sufficiently large core is a coordinate ball about
the added point, with its prescribed smooth parametrization. Under a
diffeomorphism $\widehat N\to S^4$ this is an embedded standard disk.
Shrinking the disk inside its parametrization and applying isotopy extension
identifies its complement with a standard $B^4$. Equivalently one can
straighten the diffeomorphism near the added point before removing that
point. This also identifies the prescribed end collar with a Euclidean
collar. No smooth Schoenflies statement for an arbitrary embedded $S^3$
is being assumed.
\end{proof}

\section{The minimizing property supplies an actual replacement}
Here a full-core rescaled limit means the following precise input.
Let $g_j$ be metrics on a fixed closed $M$, with
\[
 P(g_j)\longrightarrow P_*(M):=\inf_{h\in\Met(M)}P(h).
\]
For every sufficiently large fixed $R$, suppose there are embeddings
$\iota_{j,R}:N_{2R}\to M$ and numbers $\lambda_j>0$ such that
\[
 \lambda_j^{-2}\iota_{j,R}^*g_j\longrightarrow g
 \quad\hbox{in }C^2(N_{2R}).
\]
The embeddings include the whole compact core, not just an annulus or the
regular part of an orbifold.

\begin{theorem}\label{replace}
Under these assumptions, a limit satisfying the hypotheses of
Theorem \ref{standard} cannot have $0<P(g)<16\pi^2$.
\end{theorem}
\begin{proof}
By Theorem \ref{standard}, replace the metric on $N_R$ by a flat metric
whose boundary collar agrees with the Euclidean AE coordinates. On
$N_{2R}\setminus N_R$ interpolate to the pulled-back original metric,
and leave the latter unchanged near $\partial N_{2R}$.
For the limiting AE metric the transition energy is $O(R^{-2\tau})$.
For fixed $R$, $C^2$ convergence gives the same estimate up to $o_j(1)$.
The new metric is positive definite for large $R,j$, smooth, and agrees
with $g_j$ outside $\iota_{j,R}(N_{2R})$ after scaling back.
Scale invariance of $P$ gives an actual competitor $\widetilde g_j$ with
\[
 P(\widetilde g_j)
 \le P(g_j)-\int_{N_R}|\Ric_g^\circ|^2dV_g
             +C R^{-2\tau}+o_j(1).
\]
Choose $R$ first so that the removed integral exceeds $3P(g)/4$ and the
transition term is less than $P(g)/4$. Then let $j\to\infty$. This gives
\[
 \limsup_j P(\widetilde g_j)\le P_*(M)-P(g)/2,
\]
contradicting the definition of $P_*(M)$.
\end{proof}

\section{What this closes, and what it does not}
The preceding continuation, \emph{Scalar Curvature and Global Attainment}
(12 October 2026), established that a nonnegative-scalar-curvature compact
core with finite fundamental group at its connected boundary, embedded in a
closed aspherical four-manifold, has nullhomotopic inclusion. For a spherical
space-form boundary its homological consequences include $\chi(N)=1$.
Its separate ALE calculation gave
\[
 P(N)\ge4\pi^2+S(N)/12
\]
for nontrivial end group, and excluded nonflat Ricci-flat ALE cores.
The trivial end group had no positive lower bound from Gauss--Bonnet alone.

Theorem \ref{replace} supplies the missing distinction for that branch:
\emph{a small positive-energy AE core may be an admissible geometry, but
cannot survive in a genuine global minimizing sequence with full-core
smooth convergence}. For scalar-flat limits with trivial end group,
$P=0$ implies flatness: Ricci-flatness and the Euclidean asymptotic volume
ratio give the equality case of Bishop--Gromov comparison. Hence any
nonflat limit of this type has $P\ge16\pi^2$.

Combining the two results, within the stated smooth, one-ended ALE
blow-up setting on an aspherical ambient manifold, the remaining nonflat
limits have energy at least $4\pi^2$; the trivial-end branch has the stronger
$16\pi^2$ bound. Thus a sequence of total energy strictly below $4\pi^2$
has none of these limits. This is a consequence conditional on the existence
of the indicated blow-up description, not a theorem producing that
description for arbitrary minimizing sequences.

It still does not establish the desired smooth global minimizer. Nontrivial
end groups, orbifold convergence without an embedded smooth full core,
multiple-ended phenomena, and collapse have not all been excluded. In
particular, the previously analyzed scalar-flat $Q_8$ core has energy
$4\pi^2$ and is not a ball; Theorem \ref{replace} does not remove it.
The earlier doubled-hyperbolic-cusp example also shows why exclusion of
curvature bubbles alone cannot rule out long collapsing necks.

No closed non-Bach-flat global $P$ minimizer is constructed or proved to
exist in this note. The completed contribution is the standard-core and
strict-replacement argument above, with its quantified applicability.

\begin{thebibliography}{9}
\bibitem{CGY} S.-Y. A. Chang, M. J. Gursky and P. C. Yang,
\emph{A conformally invariant sphere theorem in four dimensions},
Publ. Math. IHES \textbf{98} (2003), 105--143.
\href{https://arxiv.org/abs/math/0309287}{arXiv:math/0309287}, Theorem A
and the explicit norm convention following it.
\bibitem{DM} J. Dilts and D. Maxwell,
\emph{Yamabe classification and prescribed scalar curvature in the
asymptotically Euclidean setting},
\href{https://arxiv.org/abs/1503.04172}{arXiv:1503.04172}.
\bibitem{LeBrun} C. LeBrun,
\emph{Curvature functionals, optimal metrics, and the differential topology
of 4-manifolds},
\href{https://arxiv.org/abs/math/0404251}{arXiv:math/0404251}, Section 11.
\end{thebibliography}
\end{document}
